% =============================================================================
% Section 1: Introduction
% Stubbed; theorem-first drafting per directive. Full prose pass after §3-§9 lock.
% =============================================================================

\section{Introduction}
\label{sec:intro}

\subsection{Erd\H{o}s Problem \#647 and the $\tauf$-barrier inequality}

In his 1979 paper \emph{Some unconventional problems in number theory}
\cite{Erdos1979}, P. Erd\H{o}s introduced the following problem:

\begin{problem}[Erd\H{o}s \#647]
\label{prob:e647}
Let $\tauf(n)$ denote the number of divisors of $n$. Does there exist
$n > 24$ with
\[
\max_{1 \le m < n} \bigl( m + \tauf(m) \bigr) \;\le\; n + 2 \;?
\]
\end{problem}

The case $n = 24$ holds: among $1 \le m \le 23$, the maximum of
$m + \tauf(m)$ is $25$, achieved at $m = 23$ (with $\tauf(23) = 2$),
satisfying $25 \le 26 = n + 2$. Erd\H{o}s noted that no further
example is known, and conjectured that the difference
$\max_{m<n}(m + \tauf(m)) - n$ tends to infinity.

Substituting $m = n - k$ yields the equivalent
\emph{$\tauf$-barrier inequality}:
\[
\tauf(n - k) \;\le\; k + 2 \qquad \text{for all } 1 \le k \le n - 1.
\]
Erd\H{o}s remarked of this formulation that it is ``hopeless with
our present methods,'' and that infinitely many such $n$ would
follow from Schinzel's Hypothesis~H \cite{SchinzelSierpinski1958}.

\subsection{Contribution preview}

This paper does \emph{not} resolve \cref{prob:e647}. We do not
prove that any $n > 24$ exists, nor that none does. The contribution
is structural and methodological:

\begin{enumerate}[leftmargin=*]
\item \textbf{A Lean-certified structural reduction.} Any
counterexample to \cref{prob:e647} satisfies $n \equiv r \pmod M$
where $M = 46189 = 11 \cdot 13 \cdot 17 \cdot 19$ and $r$ lies in an
explicit set of 41 residue classes (\cref{thm:stage1-reduction});
moreover every such $r$ satisfies $143 \mid r$, yielding the
canonical normalization $n = 143 \cdot Y$ with $Y \bmod 323$ in an
explicit 41-element set $\Uset \subset \Z/323\Z$
(\cref{thm:143-collapse}). Both reductions are formally verified in
Lean.

\item \textbf{A finite-local audit chain with three closure-incapable
theorems.} Within precisely-defined classes of finite-local
forced-divisor methods, no closure of the residual frontier exists at
the parameter regimes tested:
\begin{itemize}[leftmargin=*]
\item \emph{Pointwise-local-existential moving-window cover}
  (\cref{def:ple}, \cref{thm:ple-insufficiency}): the SAT instance is
  verified by \texttt{cake\_lpr} and importable into Lean via
  \texttt{Std.Tactic.BVDecide}.
\item \emph{Branch-gluing at the restricted B1 closure-capable level}
  (\cref{sec:branch-gluing-phaseB}): clean witnesses for the sentinel
  residues $r \in \{1716, 4862\}$ satisfy the necessary protected
  condition $\tau(D_k) \le k+2$ for all $k \le 1021$, the diagonal
  $\tau(D_{k_\Delta}) \le k_\Delta + 2$ at $k_\Delta = 2520\, r$, and
  the structural scan to $K_{\rm scan} = 10^7$. By
  restricted-SAT~$\Rightarrow$~unrestricted-SAT, branch-gluing is
  negative as a closure mechanism on these sentinels.
\item \emph{Composite / non-prime-power forced-divisor covers}
  (\cref{sec:composite-cover}): adversary-SAT escape on four sentinels
  $r \in \{0, 1716, 4862, 37752\}$ at $P_{\max} = 199$,
  $K_{\max} = 5000$. The encoding is mathematically a generalized
  DivisorMask at composite moduli with disjunctive cover semantics;
  every clause reduces to a valuation pattern, strengthening the
  L1--L3 DivisorMask negative result.
\end{itemize}
A diagnostic structural identity emerged from the audit
(\cref{sec:diagonal-megashift}): $F_{r, k_\Delta}(s) = A \cdot s$ at
$k_\Delta = 2520\, r$ for every residue $r$, which produces a
class-specific mega-shift kill when the chosen kernel base forces
sufficient prime divisors into $a$. Phase B shows this is
\emph{not} by itself a base-AP closure --- alternative clean kernel
bases for the sentinel residues avoid the diagonal while satisfying
the protected constraints --- but the identity unifies the
mega-shift behavior observed across all four sentinels
$r \in \{858, 1716, 2431, 4862\}$.

\item \textbf{A conditional barrier theorem under Schinzel's
Hypothesis~H.} Under H, the residual 41 classes are infinite-survivor
classes with positive Bateman--Horn singular-series density
(\cref{thm:conditional-barrier}). This makes the barrier explicit:
no \PLEclass\ finite-local forced-divisor closure of the residual
frontier exists under H, since the survivor set is genuinely infinite
at every $K$. This is consistent with Erd\H{o}s' 1979 ``hopeless with
our present methods'' assessment.

\item \textbf{An empirical floor.} No counterexample exists at the
principal residue $r = 0$ contiguously for $n \le 3.34 \times 10^{21}$,
with non-contiguous islands extending to $n \approx 1.34 \times 10^{22}$
(\cref{prop:empirical-r0}).

\item \textbf{An obstruction analysis.} The natural averaged-sieve
attacks (Selberg upper bounds, Shiu/Nair--Tenenbaum lower bounds,
multi-point $\tauf$-correlation) are blocked by four convergent
obstructions analyzed in \cref{sec:obstructions}: Selberg upper-bound
divergence, the $L_1$-vs-$L_\infty$ gap, the $k \le 2$ correlation
ceiling, and the wrong-polarity issue
\cite{TaoTeravainen2025, MRT2019, MSTT2024_I, MSTT2024_II,
Maynard2013, Maynard2014}.
\end{enumerate}

\subsection{Honest scope statement}

The 41 open-frontier residue classes remain open at the base-AP level. We
do not claim to have closed any residue. The 4 per-kernel-class
certificates produced by the structural-kill detector at
$K_{\rm scan} \le 10^7$ are class-specific and do not lift to base-AP
closures (per the quantifier audit, \cref{rem:per-kernel-class}). The
empirical floor $n \le 3.34 \times 10^{21}$ is at the principal
residue only; the 40 non-principal classes are not yet exhaustively
verified.

The contribution is best understood as a
\emph{structural barrier theorem in the spirit of the parity barrier
of Selberg \cite{Selberg1949}}: a precise
characterization of why the natural finite-local
and averaged-sieve methods cannot close \cref{prob:e647}, paired
with a concrete reduction to a small frontier and a Lean-certified
audit chain.

\subsection{Outline}

\Cref{sec:setup} fixes notation and defines the framework.
\Cref{sec:main} states the six main theorems. \Cref{sec:structural-reduction}
proves the 96~$\to$~41 Stage-1 reduction and the 143-collapse
(\cref{thm:stage1-reduction,thm:143-collapse}). \Cref{sec:csp-audit}
proves the finite-local insufficiency theorem
(\cref{thm:ple-insufficiency}) and exhibits the SAT-certified audit
chain. \Cref{sec:empirical-detail} details the empirical floor.
\Cref{sec:obstructions} analyzes the analytic obstructions.
\Cref{sec:lean} describes the Lean formalization.
\Cref{sec:open} lists open problems.

% =============================================================================
% End of section 1
% =============================================================================
